\documentclass[12pt]{article}

%\usepackage[cm]{fullpage}
\setlength{\topmargin}{-1in}
\setlength{\oddsidemargin}{-0.25in}
\setlength{\evensidemargin}{-0.25in}
\setlength{\textwidth}{7in}
\setlength{\textheight}{10in}

\usepackage{array}
\usepackage{enumitem}
\usepackage{amssymb}
\usepackage{amsmath}
\usepackage{epstopdf}
\usepackage[T1]{fontenc}
\usepackage{graphics}
\usepackage{graphicx}
\usepackage{hyperref}
\usepackage{latexsym}
\usepackage{mathrsfs}
\usepackage{multicol}
\usepackage{setspace}
\usepackage{subfigure}
\usepackage{time}
\usepackage{times}

\newcommand{\To}{\Rightarrow}
\newcommand{\del}{\nabla}
\newcommand{\R}{\mathbb{R}}
\newcommand{\C}{\mathcal{C}}
\newcommand{\N}{\mathcal{N}}
\newcommand{\A}{\mathcal{A}}
\renewcommand{\O}{\mathcal{O}}
% \renewcommand{\N}{\mathbb{N}}
\newcommand{\ben}{\begin{enumerate}}
\newcommand{\een}{\end{enumerate}}
\newcommand{\eps}{\varepsilon}
\renewcommand{\epsilon}{\varepsilon}
\newcommand{\p}{\partial}
\newcommand{\<}{\langle}
\renewcommand{\>}{\rangle}
\renewcommand{\subset}{\subseteq}
\newcommand{\cont}{\Rightarrow \Leftarrow}
\newcommand{\back}{\backslash}
\newcommand{\norm}[1]{\|{#1}\|}
\newcommand{\abs}[1]{|{#1}|}
\newcommand{\ip}[1]{\langle{#1}\rangle}
\newcommand{\m}[1]{{\mathbf{#1}}}
\newcommand{\sts}{\m{S}^\textrm{T}\m{S}}
\newcommand{\rank}{\textrm{rank}}
\newcommand{\X}{\mathcal{X}}
\newcommand{\Y}{\mathcal{Y}}
\def\ds{\displaystyle}
\def\d{\partial}
\renewcommand{\arraystretch}{1.5}
\renewcommand{\L}{\mathscr{L}}

%\includegraphics[scale = 0.65, trim = 1in 3.25in 1in 3in]{blank_axes.pdf}


\pagestyle{empty}
\begin{document}
\noindent\rule{7in}{0.4pt}\\
\begin{center}
{\sc Differential Equations\\
Order of an ODE}
\end{center}
\noindent\rule{7in}{0.4pt}
%%%%%%%%%%%%%%%%%%%%%%%%%%%%%%%%%%%%%%%%%%%%%%%%
\section*{Classifying DEs:  order}
\doublespacing
\noindent  The next idea might be easier to understand if we consider an analogous situation that has nothing
to do with differential equations.  Consider the
two equations below.\\ \vspace{0.5cm}
%	\[ x^2 - 7x  - 6 = 0 \hspace{1cm} e^{4x} = 2^{7x-8} \]

\noindent  How might you go about solving the first equation? \\ \vspace{0.5in}
%One option might be to factor the equation
%and then use properties of products that equal zero.  Another option would be to use the quadratic formula, and
%there are other viable options, as well.\\

\noindent Now consider the second equation.  Would it be appropriate to use the quadratic formula to solve it?
Why not?  The quadratic formula is great, but it does not solve every type of algebraic equation.  When we 
solve an algebraic equation, we need to inspect it and to decide which solution technique is appropriate.  
The same is true for differential equations.  We will study several different solution techniques for solving 
DEs.  Each solution technique can only be applied to certain types of DEs.  Thus, our first job is to figure
out %how to distinguish one DE from another.\\
\line(1,0){450}.\\

\noindent  The first distinguishing factor for DEs is the %{\bf order of  DE}
\line(1,0){200}, which refers to the \\ \\%highest derivative that appears in the equation.  
\line(1,0){500}.\\

\noindent {\bf Example 1:}  Consider the DE\\ \\

%	\[ x^2y'' + y''' = \sin x + y^8 \]
\noindent This DE is said to be %a third-order DE 
\line(1,0){125}\, because 
%the highest derivative that appears in the equation is a third derivative.  
\line(1,0){250}.

\noindent Do not be thrown off by the exponent 8 here that is indicating the dependent variable is to the
8th power; this has nothing to do with the order of the DE.\\

\noindent  {\bf Example 2:}  Consider the DE\\ \\

%	\[ \frac{d^2A}{dt^2} + \frac{dA}{dt} + A = 17 \]
\noindent The highest derivative is a %second derivative
\line(1,0){200}, so we say this is a \\ \\%second order DE
\line(1,0){300}.  \\
Notice that in this notation, the exponent can be used to indicate the second derivative.\\ \vspace{0.5cm}


\noindent  We have seen two different notations to indicate the derivative here. The prime notation is
quite common, but if one wants to indicate the 8th derivative using the prime notation, it gets rather
annoying.  The shorthand notation to avoid this situation is to write the dependent variable 
as what appears to be raised to a power that is in parentheses, as in %$y^{(8)}.$  
\line(1,0){100}. This notation means
the 8th derivative, not $y$ to the 8th power.\\

\end{document}